\documentclass[10pt,a4paper]{amsart}
\usepackage{amsmath,amssymb,amsthm,graphicx}
% Retained minimal preamble for the Stacks Project chapter flat.tex (More on Flatness).
% Verbatim excerpt of the chapter preamble: Xy-pic setup (native preamble lines 14-18) and the native
% theorem declarations with the mathematical macros (native preamble lines 154-217). Copied mechanically;
% nothing added, removed or redefined. The wrapper class and the amsmath/amssymb/amsthm/graphicx base
% packages are supplied by the assembler, so the class/externaldocument part of the native preamble is omitted.
% For commutative diagrams we use Xy-pic
\usepackage[all]{xy}

% We use 2cell for 2-commutative diagrams.
\xyoption{2cell}
% Theorem environments.
%
\theoremstyle{plain}
\newtheorem{theorem}[subsection]{Theorem}
\newtheorem{proposition}[subsection]{Proposition}
\newtheorem{lemma}[subsection]{Lemma}

\theoremstyle{definition}
\newtheorem{definition}[subsection]{Definition}
\newtheorem{example}[subsection]{Example}
\newtheorem{exercise}[subsection]{Exercise}
\newtheorem{situation}[subsection]{Situation}

\theoremstyle{remark}
\newtheorem{remark}[subsection]{Remark}
\newtheorem{remarks}[subsection]{Remarks}

\numberwithin{equation}{subsection}

% Macros
%
\def\lim{\mathop{\mathrm{lim}}\nolimits}
\def\colim{\mathop{\mathrm{colim}}\nolimits}
\def\Spec{\mathop{\mathrm{Spec}}}
\def\Hom{\mathop{\mathrm{Hom}}\nolimits}
\def\Ext{\mathop{\mathrm{Ext}}\nolimits}
\def\SheafHom{\mathop{\mathcal{H}\!\mathit{om}}\nolimits}
\def\SheafExt{\mathop{\mathcal{E}\!\mathit{xt}}\nolimits}
\def\Sch{\mathit{Sch}}
\def\Mor{\mathop{\mathrm{Mor}}\nolimits}
\def\Ob{\mathop{\mathrm{Ob}}\nolimits}
\def\Sh{\mathop{\mathit{Sh}}\nolimits}
\def\NL{\mathop{N\!L}\nolimits}
\def\CH{\mathop{\mathrm{CH}}\nolimits}
\def\proetale{{pro\text{-}\acute{e}tale}}
\def\etale{{\acute{e}tale}}
\def\QCoh{\mathit{QCoh}}
\def\Ker{\mathop{\mathrm{Ker}}}
\def\Im{\mathop{\mathrm{Im}}}
\def\Coker{\mathop{\mathrm{Coker}}}
\def\Coim{\mathop{\mathrm{Coim}}}

% Boxtimes
%
\DeclareMathSymbol{\boxtimes}{\mathbin}{AMSa}{"02}

%
% Macros for moduli stacks/spaces
%
\def\QCohstack{\mathcal{QC}\!\mathit{oh}}
\def\Cohstack{\mathcal{C}\!\mathit{oh}}
\def\Spacesstack{\mathcal{S}\!\mathit{paces}}
\def\Quotfunctor{\mathrm{Quot}}
\def\Hilbfunctor{\mathrm{Hilb}}
\def\Curvesstack{\mathcal{C}\!\mathit{urves}}
\def\Polarizedstack{\mathcal{P}\!\mathit{olarized}}
\def\Complexesstack{\mathcal{C}\!\mathit{omplexes}}
% \Pic is the operator that assigns to X its picard group, usage \Pic(X)
% \Picardstack_{X/B} denotes the Picard stack of X over B
% \Picardfunctor_{X/B} denotes the Picard functor of X over B
\def\Pic{\mathop{\mathrm{Pic}}\nolimits}
\def\Picardstack{\mathcal{P}\!\mathit{ic}}
\def\Picardfunctor{\mathrm{Pic}}
\def\Deformationcategory{\mathcal{D}\!\mathit{ef}}
\usepackage{cleveref}
\crefname{originaldefinition}{original definition}{original definitions}
\Crefname{originaldefinition}{Original definition}{Original definitions}
\makeatletter
\newlabel{example-two-generators}{{0EVH}{source}{}{sourcepointer-1}{}}
\newlabel{example-two-generators@cref}{{[originaldefinition][0][]0EVH}{[1][0][]source}{}{}{}}
\makeatother
\crefname{originallemma}{original lemma}{original lemmas}
\Crefname{originallemma}{Original lemma}{Original lemmas}
\makeatletter
\newlabel{lemma-h-sheaf-colim-F}{{0EVK}{source}{}{sourcepointer-2}{}}
\newlabel{lemma-h-sheaf-colim-F@cref}{{[originallemma][0][]0EVK}{[1][0][]source}{}{}{}}
\makeatother
\crefname{originallemma}{original lemma}{original lemmas}
\Crefname{originallemma}{Original lemma}{Original lemmas}
\makeatletter
\newlabel{limits-lemma-descend-invertible-modules}{{0B8W}{source}{}{sourcepointer-3}{}}
\newlabel{limits-lemma-descend-invertible-modules@cref}{{[originallemma][0][]0B8W}{[1][0][]source}{}{}{}}
\makeatother
\crefname{originallemma}{original lemma}{original lemmas}
\Crefname{originallemma}{Original lemma}{Original lemmas}
\makeatletter
\newlabel{lemma-colim-F-Vect}{{0EXB}{source}{}{sourcepointer-4}{}}
\newlabel{lemma-colim-F-Vect@cref}{{[originallemma][0][]0EXB}{[1][0][]source}{}{}{}}
\makeatother
\crefname{originallemma}{original lemma}{original lemmas}
\Crefname{originallemma}{Original lemma}{Original lemmas}
\makeatletter
\newlabel{more-morphisms-lemma-globally-generated-vanishing}{{0E7H}{source}{}{sourcepointer-5}{}}
\newlabel{more-morphisms-lemma-globally-generated-vanishing@cref}{{[originallemma][0][]0E7H}{[1][0][]source}{}{}{}}
\makeatother
\crefname{originaltheorem}{original theorem}{original theorems}
\Crefname{originaltheorem}{Original theorem}{Original theorems}
\makeatletter
\newlabel{coherent-theorem-formal-functions}{{02OC}{source}{}{sourcepointer-6}{}}
\newlabel{coherent-theorem-formal-functions@cref}{{[originaltheorem][0][]02OC}{[1][0][]source}{}{}{}}
\makeatother
\title{Frobenius-periodic vector bundles glue along the almost blowup square of an equation in two generators, giving an equivalence of categories}
\author{The Stacks Project Authors}
\date{stacks-project:flat.tex:local-snapshot-20261011; GNU FDL 1.2 or later}
\begin{document}
\maketitle
\noindent\textit{Editorial note.} All delivered mathematical text is the author's own, copied from the pinned Stacks Project chapter file flat.tex; no statement or proof step was rewritten, reordered or removed.\par\medskip
\noindent\textit{Editorial note.} The preamble is the chapter's own retained minimal preamble (sources/stacks-flat/retained-preamble.tex), a verbatim excerpt of the native Xy-pic setup and of the native theorem declarations and mathematical macros; no macro was invented or redefined, and the wrapper class and the amsmath/amssymb/amsthm/graphicx base packages are supplied by the assembler.\par\medskip
\noindent\textit{Editorial note.} The counted statement and proof contain no cross reference, no citation and no included file, so the delivered body is self contained.\par\medskip
\noindent\textit{Editorial note.} Canonical location of the counted proposition: Stacks tag 0EXD (sources/stacks-flat/pinned-original-source.tex line 12303).\par\medskip
\noindent\textit{Allowed standard prerequisites (editorial).}\begin{itemize}
\item The almost blowup square associated with an equation in two generators, and its Frobenius properties. Same chapter, flat.tex, Example labelled example-two-generators (source line 11488), invoked in the statement and at proof lines 12338 and 12346. Supplies the square X, X', Z, E and the description X = Spec(A), Z = Spec(A/(f\_1, f\_2)), X' = Proj(A[T\_0, T\_1]/(f\_2T\_0 - f\_1T\_1)), E = P\textasciicircum{}1\_Z used at proof lines 12366-12367.
\item Colimits of F-periodic vector bundles are sheaves and glue in the Zariski topology. Same chapter, Lemmas labelled lemma-colim-F-Vect (source line 12223) and lemma-h-sheaf-colim-F (source line 11855), invoked at proof lines 12346 and 12382. Supply that the colimit category of vector bundles is identified with finite locally free modules over the colimit of the structure sheaves, so that the fibre product of global section colimits computes the maps, and that objects and morphisms may be glued Zariski locally.
\item Descent of a triple of finite locally free modules along a finitely generated subalgebra, vanishing of H\textasciicircum{}1 for globally generated modules on P\textasciicircum{}1, and the theorem on formal functions. Limits, Lemma labelled limits-lemma-descend-invertible-modules (Stacks tag 0B8W) at proof lines 12368-12372; More on Morphisms, Lemma labelled more-morphisms-lemma-globally-generated-vanishing (Stacks tag 0E7H) at proof lines 12405-12406; Cohomology of Schemes, Theorem labelled coherent-theorem-formal-functions (Stacks tag 02OC) at proof lines 12420-12421. Supply the Noetherian reduction, the vanishing of the first cohomology of a globally generated module on E, and the lifting of the compatible trivialisation on E to a global map on X'; the gluing of the resulting local solutions is inside the counted span.
\end{itemize}\medskip
\ifdefined\hypertarget\hypertarget{sourcepointer-1}{}\fi
\noindent\textit{Source reference.} Original definition 0EVH: sources/stacks-flat/pinned-original-source.tex, source lines 11488--11488. Same-chapter pointer into the author environment example at source line 11488; invoked at proof line 12338 of the counted proof. It is retained only as a navigation pointer; every load-bearing step is inside the delivered span.\par\medskip
\ifdefined\hypertarget\hypertarget{sourcepointer-2}{}\fi
\noindent\textit{Source reference.} Original lemma 0EVK: sources/stacks-flat/pinned-original-source.tex, source lines 11855--11855. Same-chapter pointer into the author environment lemma at source line 11855; invoked at proof line 12346 of the counted proof. It is retained only as a navigation pointer; every load-bearing step is inside the delivered span.\par\medskip
\ifdefined\hypertarget\hypertarget{sourcepointer-3}{}\fi
\noindent\textit{Source reference.} Original lemma 0B8W: sources/stacks-chow/tags.txt, source lines 14154--14154. Cross-chapter standard prerequisite located in the project tag list staged in the local snapshot; invoked at proof line 12369 of the counted proof. It is retained only as a navigation pointer; every load-bearing step is inside the delivered span.\par\medskip
\ifdefined\hypertarget\hypertarget{sourcepointer-4}{}\fi
\noindent\textit{Source reference.} Original lemma 0EXB: sources/stacks-flat/pinned-original-source.tex, source lines 12223--12223. Same-chapter pointer into the author environment lemma at source line 12223; invoked at proof line 12382 of the counted proof. It is retained only as a navigation pointer; every load-bearing step is inside the delivered span.\par\medskip
\ifdefined\hypertarget\hypertarget{sourcepointer-5}{}\fi
\noindent\textit{Source reference.} Original lemma 0E7H: sources/stacks-chow/tags.txt, source lines 17782--17782. Cross-chapter standard prerequisite located in the project tag list staged in the local snapshot; invoked at proof line 12406 of the counted proof. It is retained only as a navigation pointer; every load-bearing step is inside the delivered span.\par\medskip
\ifdefined\hypertarget\hypertarget{sourcepointer-6}{}\fi
\noindent\textit{Source reference.} Original theorem 02OC: sources/stacks-chow/tags.txt, source lines 3511--3511. Cross-chapter standard prerequisite located in the project tag list staged in the local snapshot; invoked at proof line 12421 of the counted proof. It is retained only as a navigation pointer; every load-bearing step is inside the delivered span.\par\medskip
\begin{lemma}
\label{lemma-vector-bundle-II}
Let $p$ be a prime number. Consider an almost blowup square $X, X', Z, E$
in characteristic $p$ as in Example \ref{example-two-generators}.
Then the functor
$$
G :
\colim_F \textit{Vect}(X)
\longrightarrow
\colim_F \textit{Vect}(Z)
\times_{\colim_F \textit{Vect}(E)}
\colim_F \textit{Vect}(X')
$$
is an equivalence.
\end{lemma}

\begin{proof}
Fully faithfulness. Suppose that $(\mathcal{E}, n)$ and
$(\mathcal{F}, m)$ are objects of $\colim_F \textit{Vect}(X)$.
Let $(a, b) : G(\mathcal{E}, n) \to G(\mathcal{F}, m)$
be a morphism in the RHS. We may choose $N \gg 0$ and
think of $a$ as a map
$a : F^{N - n, *}\mathcal{E}|_Z \to F^{N - m, *}\mathcal{F}|_Z$
and $b$ as a map
$b : F^{N - n, *}\mathcal{E}|_{X'} \to F^{N - m, *}\mathcal{F}|_{X'}$
agreeing over $E$.
Choose a finite affine open covering
$X = X_1 \cup \ldots \cup X_n$ such that $\mathcal{E}|_{X_i}$
and $\mathcal{F}|_{X_i}$ are finite free $\mathcal{O}_{X_i}$-modules.
For each $i$ the base change
$$
\xymatrix{
E_i \ar[r] \ar[d] & X'_i \ar[d] \\
Z_i \ar[r] & X_i
}
$$
is another almost blow up square as in Example \ref{example-two-generators}.
For these squares we know that
$$
\colim_F H^0(X_i, \mathcal{O}_{X_i}) =
\colim_F H^0(Z_i, \mathcal{O}_{Z_i})
\times_{\colim_F H^0(E_i, \mathcal{O}_{E_i})}
\colim_F H^0(X'_i, \mathcal{O}_{X'_i})
$$
by Lemma \ref{lemma-h-sheaf-colim-F} (see proof of the lemma).
Hence after increasing $N$ we may assume
the maps $a|_{Z_i}$ and $b|_{X'_i}$ come from
maps $c_i : F^{N - n, *}\mathcal{E}|_{X_i} \to F^{N - m, *}\mathcal{F}|_{X_i}$.
After possibly increasing $N$ we may assume $c_i$ and $c_j$
agree on $X_i \cap X_j$. Thus these maps glue to give the
desired morphism $(\mathcal{E}, n) \to (\mathcal{F}, m)$
in the LHS.

\medskip\noindent
Essential surjectivity. Let $(\mathcal{F}, \mathcal{G}, \varphi)$ be a
triple consisting of
a finite locally free $\mathcal{O}_Z$-module $\mathcal{F}$,
a finite locally free $\mathcal{O}_{X'}$-module $\mathcal{G}$, and
an isomorphism $\varphi : \mathcal{F}|_E \to \mathcal{G}|_E$.
We have to show that after replacing this triple by a Frobenius
power pullback, it comes from a finite locally free $\mathcal{O}_X$-module.

\medskip\noindent
Noetherian reduction; we urge the reader to skip this paragraph.
Recall that $X = \Spec(A)$ and $Z = \Spec(A/(f_1, f_2))$,
$X' = \text{Proj}(A[T_0, T_1]/(f_2T_0 - f_1T_1))$, and
$E = \mathbf{P}^1_Z$. By Limits, Lemma
\ref{limits-lemma-descend-invertible-modules}
we can find a finitely generated $\mathbf{F}_p$-subalgebra
$A_0 \subset A$ containing $f_1$ and $f_2$ such that the triple
$(\mathcal{F}, \mathcal{G}, \varphi)$ descends to
$X_0 = \Spec(A_0)$ and $Z_0 = \Spec(A_0/(f_1, f_2))$,
$X_0' = \text{Proj}(A_0[T_0, T_1]/(f_2T_0 - f_1T_1))$, and
$E_0 = \mathbf{P}^1_{Z_0}$. Thus we may assume our schemes
are Noetherian.

\medskip\noindent
Assume $X$ is Noetherian. We may choose a finite affine open covering
$X = X_1 \cup \ldots \cup X_n$ such that $\mathcal{F}|_{Z \cap X_i}$ is free.
Since we can glue objects of $\colim_F \textit{Vect}(X)$
in the Zariski topology (Lemma \ref{lemma-colim-F-Vect}), and
since we already know
fully faithfulness over $X_i$ and $X_i \cap X_j$ (see first paragraph
of the proof), it suffices to prove the existence over each $X_i$.
This reduces us to the case discussed in the next paragraph.

\medskip\noindent
Assume $X$ is Noetherian and $\mathcal{F} = \mathcal{O}_Z^{\oplus r}$.
Using $\varphi$ we get an isomorphism
$\mathcal{O}_E^{\oplus r} \to \mathcal{G}|_E$.
Let $I = (f_1, f_2) \subset A$.
Let $\mathcal{I} \subset \mathcal{O}_{X'}$
be the ideal sheaf of $E$; it is globally generated by $f_1$ and $f_2$.
For any $n$ there is a surjection
$$
(\mathcal{I}^n/\mathcal{I}^{n + 1})^{\oplus r} =
\mathcal{I}^n/\mathcal{I}^{n + 1} \otimes_{\mathcal{O}_E}
\mathcal{G}|_E \longrightarrow
\mathcal{I}^n\mathcal{G}/\mathcal{I}^{n + 1}\mathcal{G}
$$
Hence the first cohomology group of this module is zero.
Here we use that $E = \mathbf{P}^1_Z$ and hence its structure
sheaf and in fact any globally generated quasi-coherent module
has vanishing $H^1$. Compare with More on Morphisms, Lemma
\ref{more-morphisms-lemma-globally-generated-vanishing}.
Then using the short exact sequences
$$
0 \to  \mathcal{I}^n\mathcal{G}/\mathcal{I}^{n + 1}\mathcal{G} \to
\mathcal{G}/\mathcal{I}^{n + 1}\mathcal{G} \to
\mathcal{G}/\mathcal{I}^n\mathcal{G} \to 0
$$
and induction, we see that
$$
\lim H^0(X', \mathcal{G}/\mathcal{I}^n\mathcal{G})
\to
H^0(E, \mathcal{G}|_E) = H^0(E, \mathcal{O}_E^{\oplus r}) =
A/I^{\oplus r}
$$
is surjective. By the theorem on formal functions
(Cohomology of Schemes, Theorem \ref{coherent-theorem-formal-functions})
this implies that
$$
H^0(X', \mathcal{G}) \to
H^0(E, \mathcal{G}|_E) = H^0(E, \mathcal{O}_E^{\oplus r}) =
A/I^{\oplus r}
$$
is surjective. Thus we can choose a map
$\alpha : \mathcal{O}_{X'}^{\oplus r} \to \mathcal{G}$
which is compatible with the given trivialization
of $\mathcal{G}|_E$. Thus $\alpha$ is an isomorphism over
an open neighbourhood of $E$ in $X'$. Thus every point
of $Z$ has an affine open neighbourhood where
we can solve the problem. Since $X' \setminus E \to X \setminus Z$
is an isomorphism, the same holds for points of $X$ not in $Z$.
Thus another Zariski glueing argument finishes the proof.
\end{proof}

\end{document}
